\newpage
\textbf{\huge Chapter 2}
\section{LITERATURE REVIEW}
\begin{justify}
This chapter reviews prior work in three areas relevant to this project: sentiment analysis (lexicon-based through neural), topic modelling, and studies focused on Amazon product review data. Understanding where existing approaches fall short was important for justifying the design decisions made here.
\end{justify}

\subsection{Sentiment Analysis}

\subsubsection{Lexicon-Based Methods}
\begin{justify}
The earliest sentiment analysis work relied on manually built lexicons. SentiWordNet \cite{esuli2006sentiwordnet} extended WordNet by assigning positive, negative, and objective scores to word senses. Hu and Liu \cite{hu2004mining} applied a lexicon-based approach to a small set of Amazon product pages to perform aspect-level sentiment analysis using adjective polarity — one of the earlier examples of domain-specific review analysis.

VADER (Valence Aware Dictionary and sEntiment Reasoner) \cite{vader}, developed by Hutto and Gilbert in 2014, remains probably the most widely used lexicon-based tool. It was designed specifically for social media text and uses five hand-coded heuristics on top of its 7,517-word lexicon to handle things like punctuation emphasis, ALL CAPS, degree modifiers, and contractions. The compound score is computed as:

\begin{equation}
    s_c = \frac{\sum s_i}{\sqrt{\sum s_i^2 + \alpha}}
\end{equation}

where $s_i$ are token-level valence scores and $\alpha = 15$ normalises the range to $[-1, 1]$. VADER works well on short informal text, but it has no notion of sentence-level context — it scores tokens independently, which causes problems on product reviews where sentiment often builds across multiple clauses.
\end{justify}

\subsubsection{Machine Learning Approaches}
\begin{justify}
Before transformers took over, classical ML methods with bag-of-words features achieved strong results. Pang et al.\ \cite{pang2002thumbs} showed that SVMs with unigram features outperformed Naïve Bayes on movie review classification. Zhang et al.\ \cite{zhang2015character} pushed accuracy to 95.7\% on Amazon reviews using character-level CNNs. Maas et al.\ \cite{maas2011learning} incorporated sentiment information directly into word vectors for binary classification.

The shared limitation of all these approaches is that they operate on surface-level features. Long-range dependencies, sarcasm, and contextual negation are out of reach without sequence-level modelling.
\end{justify}

\subsubsection{Transformer-Based Models}
\begin{justify}
BERT \cite{devlin2019bert} changed the field significantly. Trained on 3.3 billion words using masked language modelling and next-sentence prediction, its bidirectional attention mechanism captures context from both directions simultaneously. Fine-tuning BERT on the Stanford Sentiment Treebank (SST-2) \cite{socher2013recursive} became the standard recipe for sentiment classification for a while.

DistilBERT \cite{sanh2019distilbert} — a knowledge-distilled version of BERT that is 40\% smaller and 60\% faster while retaining 97\% of BERT's performance — is the model used here as a transformer baseline. However, both BERT and DistilBERT show a noticeable domain mismatch on product reviews: fine-tuned on polished movie criticism, they misclassify casual e-commerce language as negative when no strong positive signal is present.

Liu et al.\ \cite{liu2019roberta} addressed some of BERT's training weaknesses with RoBERTa, removing next-sentence prediction, using dynamic masking, larger batches, and training on 160GB of data. RoBERTa consistently outperforms BERT on benchmarks. Building on this, Barbieri et al.\ \cite{barbieri2020tweeteval} released TweetEval and fine-tuned RoBERTa on 124 million tweets, producing \texttt{cardiffnlp/twitter-roberta-base-sentiment} — a three-class (positive, neutral, negative) classifier naturally suited to informal, short-form text.

This is the model selected as the primary sentiment classifier in this project. Yin et al.\ \cite{yin2020benchmarking} benchmarked 17 models on Amazon reviews and found that tweet-trained transformers outperformed movie-review-trained ones by 4--8 percentage points — a finding that directly informed this choice.
\end{justify}

\subsubsection{Implicit Sentiment}
\begin{justify}
One underappreciated problem in review sentiment is implicit negativity. Zhao et al.\ \cite{zhao2023implicit} categorise this into hedging, rhetorical questions, presupposition, and sarcasm, noting that standard classifiers miss 15--25\% of negative reviews that use these constructions. Phrases like ``not exactly what I expected'' or ``I guess it works'' are clearly dissatisfied but score near-neutral under both lexicon and transformer approaches. The implicit negativity detection component in this pipeline draws on this line of work, using regex patterns to flag these constructions before fusion \cite{hancock2007deception}.
\end{justify}

\subsection{Topic Modelling}

\subsubsection{Latent Dirichlet Allocation}
\begin{justify}
LDA \cite{blei2003latent} by Blei, Ng, and Jordan (2003) remains the textbook probabilistic topic model. It treats documents as mixtures of topics and topics as distributions over words. The generative process for document $d$ is:

\begin{align}
    \theta_d &\sim \text{Dir}(\alpha) \\
    \phi_k   &\sim \text{Dir}(\beta) \\
    z_{dn}   &\sim \text{Multinomial}(\theta_d) \\
    w_{dn}   &\sim \text{Multinomial}(\phi_{z_{dn}})
\end{align}

where $\theta_d$ is the per-document topic distribution, $\phi_k$ is the per-topic word distribution, and $z_{dn}$ is the topic assigned to the $n$-th word of document $d$. Inference is done via variational EM or Gibbs sampling. LDA has two well-known issues for this kind of work: you have to specify the number of topics $K$ in advance, and it uses bag-of-words representations that throw away word order and semantics.
\end{justify}

\subsubsection{Neural Topic Models}
\begin{justify}
Neural topic models tried to address LDA's bag-of-words limitation by incorporating word embeddings. Top2Vec \cite{angelov2020top2vec} jointly embeds documents and words and clusters them in the same semantic space. The main drawback is that these models need to be retrained from scratch on any new corpus, and they lack the modularity that makes pipeline-based approaches easier to maintain and update.
\end{justify}

\subsubsection{BERTopic}
\begin{justify}
BERTopic \cite{grootendorst2022bertopic}, introduced by Grootendorst in 2022, is currently the state-of-the-art approach for short-text topic modelling. It works in three main stages:

\begin{enumerate}
    \item \textbf{Embedding}: Reviews are encoded with a pre-trained sentence transformer (\texttt{all-MiniLM-L6-v2} \cite{reimers2019sentencebert}), producing 384-dimensional vectors $\mathbf{e}_d \in \mathbb{R}^{384}$ optimised for semantic similarity.

    \item \textbf{Dimensionality Reduction}: UMAP \cite{mcinnes2018umap} compresses these down to 5 dimensions while preserving the local and global structure of the original space:
    \begin{equation}
        \min_{\mathbf{y}} \sum_{ij} \left[ p_{ij} \log \frac{p_{ij}}{q_{ij}} + (1-p_{ij})\log\frac{1-p_{ij}}{1-q_{ij}} \right]
    \end{equation}
    where $p_{ij}$ and $q_{ij}$ are fuzzy set memberships in the high- and low-dimensional spaces.

    \item \textbf{Clustering}: HDBSCAN \cite{campello2013density} groups the reduced embeddings into variable-density clusters without needing the number of topics specified upfront. Documents in sparse regions get labelled $-1$ (noise).

    \item \textbf{Topic Representation}: For each cluster, c-TF-IDF identifies the most representative terms:
    \begin{equation}
        \text{c-TF-IDF}_{t,c} = \text{tf}_{t,c} \cdot \log\left(1 + \frac{A}{f_t}\right)
    \end{equation}
    where $\text{tf}_{t,c}$ is term frequency in class $c$, $A$ is the average words per class, and $f_t$ is the total count of the term. The top-4 terms form the topic label.
\end{enumerate}

Grootendorst shows that BERTopic achieves higher NPMI coherence than LDA and Top2Vec on short-text corpora — exactly the setting of product reviews.
\end{justify}

\subsection{Amazon Review Analytics}
\begin{justify}
McAuley and Leskovec \cite{mcauley2013hidden} studied the relationship between review text and star ratings in the Amazon dataset, finding that text carries substantially more signal than numeric ratings. He and McAuley \cite{he2016ups} later released the multi-category Amazon review dataset that is widely used in research, noting meaningful sentiment distribution differences across product categories.

Pontiki et al.\ \cite{pontiki2014semeval} formalised Aspect-Based Sentiment Analysis (ABSA) through the SemEval challenge series, decomposing sentiment into aspect-opinion pairs. More recent work \cite{xu2020bert} applies BERT to ABSA with strong results. The aspect sentiment component in this project takes inspiration from this line of work but uses a keyword-matching approach instead of span extraction, which trades off precision for scalability over 172,000 reviews.

Table~\ref{tab:literature_comparison} summarises the key related works and how this project relates to them.
\end{justify}

\begin{table}[H]
\centering
\small
\caption{Comparison of Related Work in Review Sentiment and Topic Analysis}
\renewcommand{\arraystretch}{1.3}
\resizebox{\textwidth}{!}{
\begin{tabular}{|p{2.5cm}|p{3cm}|p{3cm}|p{3cm}|p{3cm}|}
\hline
\textbf{Reference} & \textbf{Approach} & \textbf{Dataset} & \textbf{Task} & \textbf{Limitation} \\
\hline
Hu \& Liu \cite{hu2004mining} & Lexicon + rules & Amazon (5 products) & Aspect sentiment & Small scale, domain-specific \\
\hline
Blei et al.\ \cite{blei2003latent} & LDA & News/Wikipedia & Topic modelling & Bag-of-words, K required \\
\hline
Maas et al.\ \cite{maas2011learning} & Word vectors & IMDB & Sentiment & Binary, domain-specific \\
\hline
Devlin et al.\ \cite{devlin2019bert} & BERT fine-tune & SST-2, GLUE & Sentiment, NLU & Domain mismatch for reviews \\
\hline
Barbieri et al.\ \cite{barbieri2020tweeteval} & RoBERTa tweets & TweetEval & 3-class sentiment & Short text, tweet domain \\
\hline
Grootendorst \cite{grootendorst2022bertopic} & BERTopic & ArXiv, 20News & Topic modelling & High noise for short text \\
\hline
\textbf{This work} & VADER + RoBERTa + BERTopic + BI & 172K Amazon reviews & Sentiment + Topics + BI & -- \\
\hline
\end{tabular}
}
\label{tab:literature_comparison}
\end{table}
